\documentclass[12pt,a4paper,onecolumn]{exam}
\usepackage{amsmath}
\usepackage{amssymb}
\usepackage{graphicx}
\usepackage{float}
\usepackage{booktabs}
\usepackage{geometry}
\usepackage{tikz}
\usepackage[skins]{tcolorbox}
\geometry{margin=1.8cm}


% --- Define our simple \questionheader command ---
\newcommand{\questionheader}[1]{%
  \begin{tcolorbox}[
    enhanced,
    colback=black,
    coltext=white,
    boxrule=0pt,
    fontupper=\Large\bfseries,
    arc=4mm
  ]
  #1
  \end{tcolorbox}%
}
% --- End of definition ---
\newtcolorbox{answerbox}[1]{
  boxrule=0.4pt,
  colback=white,
  height=#1,
}


\usepackage[font = small]{caption}
\usepackage{subcaption}

\newenvironment{blanksolution}
  {%
    \renewcommand{\solutiontitle}{\noindent}%
    \begin{solution}%
  }%
  {\end{solution}}
\usepackage{listings}

\lstset{
    language=Python,
    basicstyle=\ttfamily,
    }

\pagestyle{headandfoot}
\runningheader{Assignment 9}{Signal Processing in Practise}{Dwaipayan Haldar}

\begin{document}

\begingroup
    \centering
    \LARGE E9 222 Signal Processing in Practice\\
    \LARGE Assignment 9\\[0.5em]
    \large \today\\[0.5em]
    \large Dwaipayan Haldar \par
\endgroup
\noindent\rule{\textwidth}{0.5pt}
\printanswers
\renewcommand{\solutiontitle}{\noindent\textbf{Ans:}\enspace}
\setlength{\parskip}{2pt}


%% ─────────────────────────────────────────────────────────────────────────────
\questionheader{Question 1: Supervised Training}

\begin{solution}

\textbf{Setup.}
ResNet-9 was trained from scratch on the 300 labeled images (100 per class: cat, deer, dog)
for 200 epochs using cross-entropy loss.
Optimizer: Adam ($\text{lr}=10^{-3}$, weight decay $10^{-4}$) with cosine-annealing LR schedule.
Batch size: 32.
The checkpoint with the best validation accuracy was retained.
Weak augmentation (random horizontal and vertical flip) was applied to labeled samples during training.
No unlabeled data was used.

\textbf{Results.}
The final test accuracy was \textbf{57.50\%} (best val accuracy: 63.33\%).

\begin{figure}[H]
    \centering
    \includegraphics[width=\linewidth]{q1_supervised_results.png}
    \caption{Q1 supervised baseline. \textit{Left:} Training CE loss decreases steadily and plateaus
    around epoch 100, indicating convergence with no further improvement. \textit{Centre:} Train
    accuracy reaches $\sim$100\% while validation accuracy saturates around 60--63\%, a clear sign of
    overfitting due to the very small labeled set (300 samples). \textit{Right:} The confusion matrix
    shows the model handles \textit{deer} well but frequently confuses \textit{cat} and \textit{dog},
    which share visual texture similarities at 32$\times$32 resolution.}
    \label{fig:q1}
\end{figure}

\end{solution}


%% ─────────────────────────────────────────────────────────────────────────────
\questionheader{Question 2: FixMatch}

\begin{solution}

\textbf{Setup.}
FixMatch uses both the 300 labeled images and 5\,400 unlabeled images.
Labeled batch size $B=32$; unlabeled batch size $\mu B = 128$ ($\mu=4$).
Per epoch the unlabeled loader is iterated once while the labeled loader is cycled.
The total loss is $\mathcal{L} = \mathcal{L}_s + \lambda\,\mathcal{L}_u$, where
$\mathcal{L}_u$ is computed only for unlabeled samples whose weak-view confidence exceeds threshold $\tau$.

\vspace{4pt}
\noindent\textbf{Augmentations.}
\begin{itemize}\setlength\itemsep{0pt}
  \item \textbf{Weak}: \texttt{RandomHorizontalFlip} + \texttt{RandomVerticalFlip}.
  \item \textbf{Strong}: same flips + \texttt{RandAugment(num\_ops=2, magnitude=9)} +
        \texttt{RandomErasing(p=0.5)} (applied post-\texttt{ToTensor}).
\end{itemize}


\vspace{4pt}
\noindent\textbf{Hyperparameter sweep.}
A grid search was run over $\tau \in \{0.70, 0.80, 0.90, 0.95\}$ and
$\lambda \in \{0.5, 1.0, 2.0\}$ for 100 epochs each (12 runs total). The pseudocode in the assignment uses SGD but I used Adam because of better convergence. Also, I have plotted pseudo label utilization, that gives a good idea how certain the model is, as the epoch increases the utilization also increases. 

\begin{figure}[H]
    \centering
    \includegraphics[width=0.55\linewidth]{fixmatch_sweep_heatmap.png}
    \caption{Val accuracy heatmap across the hyperparameter sweep.
    The red star marks the best configuration: $\tau=0.90$, $\lambda=0.5$ (val acc 72.0\%).
    Higher thresholds ($\tau=0.95$) restrict pseudo-labels too aggressively and slightly
    hurt performance. A smaller $\lambda=0.5$ works best at $\tau=0.9$ because the
    pseudo-labels, although confident, are not yet perfectly reliable; a lower weight
    prevents the unsupervised loss from overwhelming the supervised signal.}
    \label{fig:sweep}
\end{figure}

\noindent\textbf{Full training with best config ($\tau=0.90$, $\lambda=0.5$, 200 epochs).}
Test accuracy: \textbf{71.37\%}.

\begin{figure}[H]
    \centering
    \includegraphics[width=\linewidth]{q2_fixmatch_results.png}
    \caption{FixMatch full run ($\tau=0.90$, $\lambda=0.5$).
    \textit{Top-left:} The total loss is dominated by $\mathcal{L}_s$ early on; the
    unsupervised component $\mathcal{L}_u$ rises gradually as more pseudo-labels pass
    the threshold, then stabilises.
    \textit{Top-right:} Mask ratio starts below 0.4 and climbs to $\sim$0.8 by
    epoch 50, confirming that the supervised bootstrapping phase unlocks the unlabeled
    data progressively.
    \textit{Bottom-left:} FixMatch val accuracy (orange) overtakes the supervised
    baseline (blue) from around epoch 20 and continues improving, demonstrating that
    unlabeled data provides a consistent regularisation benefit.
    \textit{Bottom-right:} The confusion matrix is more balanced than the baseline;
    cat--dog confusion is visibly reduced.}
    \label{fig:q2}
\end{figure}

\end{solution}

\questionheader{Question 3: Analysis and Observation}

\begin{solution}

\begin{table}[H]
\centering
\small
\begin{tabular}{lcc}
\toprule
\textbf{Method} & \textbf{Val Acc} & \textbf{Test Acc} \\
\midrule
Supervised baseline (Q1) & 63.33\% & 57.50\% \\
FixMatch $\tau=0.90$, $\lambda=0.5$ (Q2) & 74.67\% & \textbf{71.37\%} \\
\bottomrule
\end{tabular}
\caption{Summary of results. FixMatch improves test accuracy by \textbf{+13.87 percentage points}.}
\end{table}

\begin{figure}[H]
    \centering
    \begin{subfigure}[t]{0.48\linewidth}
        \includegraphics[width=\linewidth]{q3_val_accuracy_comparison.png}
        \caption{Validation accuracy over 200 epochs for both methods (dotted lines show final test accuracy).}
    \end{subfigure}
    \hfill
    \begin{subfigure}[t]{0.48\linewidth}
        \includegraphics[width=\linewidth]{q3_comparison_cm.png}
        \caption{Side-by-side test confusion matrices.}
    \end{subfigure}
    \caption{Q3 comparison: supervised baseline (blue) vs.\ FixMatch (orange).}
    \label{fig:q3}
\end{figure}

\noindent\textbf{Observations.}
The val-accuracy plot (Figure~\ref{fig:q3}a) shows that the supervised baseline plateaus
sharply around epoch 30--40 and never recovers, a direct consequence of overfitting to
only 300 labeled samples.
FixMatch, by contrast, continues to improve well beyond epoch 100 as rising mask
utilisation allows the 5\,400 unlabeled images to act as implicit regularisers.
The gap between the dotted test lines ($+13.87$\,pp) confirms the benefit is not
just on the validation set but generalises.

The confusion matrices (Figure~\ref{fig:q3}b) reveal that the gain is concentrated
on the cat and dog classes: the supervised baseline confuses 257 out of 1\,000 cats
as dogs (and vice versa), whereas FixMatch reduces these cross-class errors substantially,
suggesting that consistency regularisation on the abundant unlabeled images sharpens the
feature boundary between visually similar fine-grained classes.

The sweep (Figure~\ref{fig:sweep}) shows that the sweet spot is $\tau=0.90$:
too low ($\tau=0.70$) admits noisy pseudo-labels and the model learns from incorrect
supervision, while too high ($\tau=0.95$) leaves most unlabeled data unused,
particularly in early epochs when confidence is low.
A moderate unlabeled weight $\lambda=0.5$ performs best at the optimal threshold,
balancing the reliable supervised signal against the inherently noisier pseudo-label signal.

\end{solution}
\end{document}
